
\subsubsection{UTS description}
The Unbalanced Tree Search benchmark measures the rate of traversal of
a tree generated on the fly using a splittable random number
generator. For this submission, we used a \emph{fixed geometric law}
for the random number generator with branching factor $b_0=4$, seed
$r=19$, and tree depth $d$ varying from 13 to 20 depending on core counts and architecture.

\begin{quote}
The nodes in a geometric tree have a branching factor that follows a
geometric distribution with an expected value that is specified by the  
parameter $b_0 > 1$. The parameter $d$ specifies its maximum depth
cut-off, beyond which the tree is not allowed to grow ... The expected
size of these trees is $(b_0)^d$, but since the geometric distribution
has a long tail, some nodes will have significantly more than $b_0$
children, yielding unbalanced trees. 
\end{quote}

The depth cut-off makes it possible to estimate the size of the trees
and shoot for a target execution time. To be fair, the distance from
the root to a particular node is never used in our benchmark
implementation to predict the size of a subtree. In other words, all
nodes are treated equally, irrespective of the current depth.  

\subsubsection{Design}

The material in this section is excerpted from \cite{glb}, and
unchanged from our 2012 HPCC submission.

One common way to load balance is to use {\em work-stealing}. For
shared memory system this technique has been pioneered in the Cilk
system.  Each worker maintains a double ended queue (deque) of tasks. When a
worker encounters a new task, it pushes its continuation onto the
bottom of the deque, and descends into the task. On completion of the task, it
checks to see if its deque has any work, if so it pops a task (from the
bottom) and continues with the work. Else (its deque is empty) it
looks for work on other workers' deques. It randomly determines a {\em
  victim} worker, checks if its queue is non-empty, and if so, pops a
task from the top of the deque (the end other than the one being
operated on by the owner of the deque).  If the queue is empty, it
guesses again and continues until it finds work. Work-stealing systems
are optimized for the case in which the number of steals is much
smaller than the total amount of work. This is called the ``work
first'' principle -- the work of performing a steal is incurred by the
thief (which has no work), and not the victim (which is busy
performing its work). 

Distributed work-stealing must deal with some additional
complications. First, the cost of stealing is usually significantly
higher than in the shared memory case. For many systems, the target
CPU must be involved in processing a steal attempt. Additionally, one
must solve the distributed termination detection problem. The system
must detect when all workers have finished their work, and there is no
more work. At this point all workers should be shut down, and the
computation should terminate. 

X10 already offers a mechanism for distributed termination detection
-- \texttt{finish}. Thus, in principle it should be possible to spawn
an activity at each place, and let each look for work. To trigger
\texttt{finish} based termination detection however, these workers
must eventually terminate. But when? One simple idea is that a worker
should terminate after $k$ steal attempts have failed. However this
leaves open the possibility that a worker may terminate too early --
just because it happened to be unlucky in its first $k$ guesses. If
there is work somewhere else in the network then this work
can no longer be shared with these terminated workers, thereby
affecting parallel efficiency. (In an extreme case this could lead to
sequentialize significant portion of the work.) 

Therefore there must be a way by which a new activity can be launched 
at a place whose worker has already terminated. This leads to the idea
of a {\em lifeline graph}. For each place $p$ we pre-determine a set of
other places, called {\em buddies}. Once the worker at $p$ has
performed $k$ successive (unsuccessful) steals, it examines its
buddies in turn. At each buddy it determines whether there is some
work, and if so, steals a portion. However, if there is no work, it
records at the buddy that $p$ is waiting for work. If $p$ cannot find
any work after it has examined all its buddies, it dies -- the place $p$
now becomes quiescent.  

But if $P$ went to a buddy $B$, and $B$ did not have any work, then
it must itself be out looking for work -- hence it is possible that it
will soon acquire work. In this case we require that $B$ {\em
  distribute} a portion of the work to those workers that attempted to
buddy steal from it but failed. Work must be spawned on its own
\texttt{async} -- using the \texttt{at (p) async S} idiom. If $P$ had
 no activity left, it now will have a fresh activity. Thus, unlike
 pure work-stealing based systems, a system with lifeline graphs will
 see its nodes moving from a state of processing work (the active state), to a state in
 which they are stealing to a state in which they are dead, to a state
 (optionally) in which they are woken up again with more work (and are
 hence back in the active state).  

Note that when there is no work in the system, all places will be dead,
there will be no active \texttt{async} in the system and hence the
top-level \texttt{finish} can terminate. 

The only question left is to determine the assignment of buddies to a
place. We are looking for a directed graph that is fully connected (so
work can flow from any vertex to any other vertex) and that has a low
diameter (so latency for work distribution is low) and has a low
degree (so the number of buddies potentially sending work to a dead
vertex is low).  $z$-dimensional hyper-cubes  satisfy these properties
and have been implemented. 

\subsubsection{Implementation}\label{sec:uts-inline}
Our implementation follows from \cite{glb} and includes the
performance improvements described in our HPCC 2012 submission
(using \code{FINISH\_DENSE} to reduce the overhead of termination
detection), using \code{@Uncounted} asyncs,  and using a compact
representation of the active nodes in a place using an implicit (lower,
upper) bounds representation rather than an explicit enumeration.

We have made a change in how the native SHA1 code is integrated into
X10. Consequently we replaced all of the 169 lines of very complex C++ wrapper code
with 3 native annotations (i.e., 3 lines of C code embedded into X10) as for instance in:

{\small
\begin{verbatim}
def this(seed:Int, depth:Int) {
  @Native("c++", "rng_init((RNG_state*)this, seed);") {}
  this.depth = depth;
}
\end{verbatim}
}

We fixed a nasty deadlock condition when using multi-dimensional lifeline graphs.


\subsubsection{Verification}
The code was validated by checking the results for specific
configurations against pre-computed numbers. 


% Runs on MPI now ... were there any issues here?
\subsubsection{UTS-GLB}\label{sec:utsg}
UTS-G is a variant of UTS written using the GLB library. The code for
\code{GlobalLoadBalancer} and \code{LocalJobRunner} is no longer
needed, since they are part of the library. The code for managing bags
of tasks is refactored into a \code{UTSTaskBag} class. Similarly the
code for executing a UTS task \code{runATMostNTasks} in
\code{UTSTaskFrame} is similarly very straightforward.  The native
code (for SHA1 hashes) remains unchanged -- it belongs in the
application and not the library.

Thanks to the GLB library the number of X10 LOCs drops from 493 for UTS to 225 for UTS-G.

%% \subsection{Performance}
%% \begin{figure}
%% \centering
%% \includegraphics[width=3.25in]{figures/p7ih/uts-p7ih.eps}
%% \caption{X10 Performance for UTS (on Power~775)}
%% \label{fig:perf:uts-p7ih}

%% \includegraphics[width=3.25in]{figures/bgq/uts-bgq.eps}
%% \caption{X10 Performance for UTS (on Blue~Gene/Q)}
%% \label{fig:perf:uts-bgq}

%% \includegraphics[width=3.25in]{figures/k/uts-k.eps}
%% \caption{X10 Performance for UTS (on K)}
%% \label{fig:perf:uts-k}
%% \end{figure}

%% We report performance for UTS and UTS-G.

%% On Power~775, as we demonstrated in our 2012 HPCC Submission, UTS achieves linear speedup up to 50,000 cores. UTS-G is only marginally slower.

%% On Blue~Gene/Q and K however, we ran the code ``as is'' without having the time to tune the main parameters of the load balancer to the specifics of the architecture. These parameters include the mapping of the lifeline graph onto the network axes (especially the dimensionality) as well as the frequency of network probes. With proper tuning, we believe we can redress the curves dramatically. In essence, with our implementation and our 2012 runs at scale on Power~775, we established that the benchmark (for this family of trees at least) is compute bound.

%% On Blue~Gene/Q, the sequential performance of the UTS-G code is 15\% less than the UTS code: the higher-order implementation confuses the optimizer of the backend compiler (xlC). This is not the case on other platforms. As a result, UTS-G appears to scale better than UTS, but this is misleading.
%% At scale, they are more or less on par as UTS only achieves 60\% efficiency for now. Hence the sequential inefficiency of UTS-G can hide behind the latency of the load balancer.

%% The measurements on K are very preliminary. The code runs and computes correct results. The performance is under investigation.



%We discuss the performance together with UTS-G in Section~\ref{sec:utsg}.
%\todo{Performance analysis. Why are numbers bad for the K. Is
%  there some basic operation that is slow on the K?}
